% Admitted by research/round28/advisor/ai4-gate.json after final skeptical review.
\section{HNM AI4: a uniform family of complete physical ratio tests\workstar}
AI3 identified a remaining spatial certificate at small $u$. AI4 adaptively shortens the prescribed duration to $z=10^{-7}$ and increases the proof collar with a dyadic schedule. Under the inherited canonical full-link transfer theorems, the two complete predicted ratio intervals have gap at least $39u/40$ for \emph{every real} $0<u\le10^{-24}$. Both independent producers prove this common conclusion using different explicit upper bounds. Every state, spatial and averaging error remains; this is a conditional family of candidate-pair tests, not a continuous inverse or an apparatus.

\subsection{The family and its unchanged physical clock}
Fix $u_0=10^{-24}$, $z=10^{-7}$ and let $j\ge0$ be the unique integer for which
\begin{equation}
 \frac{U_j}{2}<u\le U_j,\qquad U_j=u_0,2^{-j},\qquad k=6+j.
 \label{r28:eq:ai4-bands}
\end{equation}
An excluded lower endpoint belongs to the next band. The candidates remain \eqref{r28:eq:ai3-candidates}; hence $\eta(1-q)^3=u^3/2$ in both models and their common physical time is
\begin{equation}
 t=2zu^{-3}\frac{\hbar}{\alpha}.
 \label{r28:eq:ai4-clock}
\end{equation}
At the base upper endpoint it is already $2\times10^{65}\hbar/\alpha$ and grows without bound. Positive $\alpha$, $\hbar$, $\alpha/E_*$ and spacing remain fixed. Changing $z$ changes the observation duration within the admitted range, not the Hamiltonian or a separately fitted candidate clock. The values $q,\eta$ are canonical-model parameters, not lattice spacing.

The actual probes, multiplication norm $2\sqrt2$, common ten-factor seed, reached sixteen-cycle component and exterior-loading hypotheses are those of AI3. The complete physical radius remains \eqref{r28:eq:ai3-physical-radius}, now at the shorter duration and the scheduled $k$. Both connected-state means, both evolution columns, complete factors, and every omitted spatial word remain. The reached averaged component is not the whole electric shell or an instantaneous invariant space. The following proof uses source-bound AA2/AD2/Y1/Y2/AI3 transfer premises; it does not replace them with matrices or revalidate the entire historical proof chain.

\subsection{A collar envelope proved at every depth}
A unit face touching a current link adds at most one to each endpoint-box coordinate. Completing a strip owner adds at most three in $x$, one in $y$ and zero in $z$; its actual six $x$ links and four $y$ links are all included. A free factor adds no completion. Induction on the complete owner rule therefore gives
\begin{equation}
 \operatorname{vertices}(E(F_k))\subset
 [0,4(k+1)]\times[0,2(k+1)]\times[0,k+1]
 \quad\text{for every }k\ge0.
 \label{r28:eq:ai4-geometry}
\end{equation}
Negative-coordinate faces are absent. Every retained face lies wholly inside this box. Counting unit faces in its three orientations yields the forward envelope, with $L=k+1$,
\begin{equation}
 N_k\le\overline N_k=24L^3+14L^2,\qquad
 M_k\le\overline M_k=L^2(L+7/12).
 \label{r28:eq:ai4-forward-envelope}
\end{equation}
The reverse proof instead counts three orientations at every possible anchor, obtaining the valid larger bound
\begin{equation}
 N_k\le3(4k+5)(2k+3)(k+2),\qquad
 M_{6+j}\le435(j+1)^3.
 \label{r28:eq:ai4-reverse-envelope}
\end{equation}
At depth six, the exact admitted face count is 1808, while the two envelopes are 8918 and 10440. Every admitted depth-zero-through-six complete link is checked against the new geometry. No actual collar beyond depth six is enumerated; a large-depth polynomial envelope is never presented as an exact face count.

Hereafter the displayed detailed proof follows the forward constants. Put
\begin{equation}
 M_j=(j+7)^2(j+7+7/12),\qquad M_0=4459/12,
 \qquad x=15z.
\end{equation}
The full spatial tail in \eqref{r28:eq:ai3-collar-tail} satisfies
\begin{equation}
 \frac{D_{k+1}}{D_k}=\frac{x}{1-x}\frac{k+13/3}{k+2},\qquad
 \frac{D_{6+j}}u\le\frac{2^{j+1}D_{6+j}}{u_0}
 \le\frac{2D_6}{u_0},
 \label{r28:eq:ai4-spatial-induction}
\end{equation}
because the successive ratio of the middle majorants is at most
$2x(31/24)/(1-x)<1$. This proves every integer $j$, rather than extrapolating from sampled bands.

Each of the three positive affine factors in $M_j$ grows by at most $8/7$ per step. Thus $M_{j+1}\le gM_j$, with $g=(8/7)^3$, and
\begin{equation}
 \frac{g^2}{4}<1,\qquad\frac{g^2}{8}<1.
 \label{r28:eq:ai4-growing-averaging}
\end{equation}
Consequently both $U_j^2M_j(1+3zM_j)$ and $U_j^3M_j(1+3zM_j)$ decrease. The full quadratic $M_j^2$ term is retained. Since both candidates obey $\tau_q\le3u^3/4$, the averaging contribution is bounded by $48u^3M_j(1+3zM_j)$, with its normalized and absolute maxima controlled at the base band. The reverse proof independently bounds the cubic and sixth-power sequences $(j+1)^3/4^j\le2$ and $(j+1)^6/4^j\le64$ by finite algebraic prefixes followed by strict ratio induction. Those prefixes are not additional physical fixtures.

\subsection{Continuous bounds and the location of arithmetic error}
For the entire family $q\ge q_*=99/100$. With $e=1-q$, positive polynomial inequalities for the unchanged $p(q),b(q)$ give
\begin{equation}
 \frac z{100}\le v\le V:=\frac z{64},\qquad
 \tau_q\le\frac32\eta e^3,\qquad
 \sqrt{b(q^2)/96}\le\frac1{80e^{3/2}}.
 \label{r28:eq:ai4-continuous}
\end{equation}
For example $p(q)-2(1+q)^2(1+q^2)=q+q^2+2q^3+q^4\ge0$ gives the upper $v$ bound; rational endpoint inequalities at $q_*$ give the others. No grid interpolation supplies these continuous claims.

The complete forward state bounds for $H_1,H_2$ are $(36/5)u^{3/2}$ and $(36/25)u^{3/2}$. An absolute mesh $\delta=2^{-320}$ applied to the \emph{inner} square root adds at most
\begin{equation}
 \frac{384\tau_q\delta}{1-\eta},\qquad
 \text{hence at most }576\delta u^3\text{ or }(1536/5)\delta u^3.
 \label{r28:eq:ai4-rounding}
\end{equation}
The prefactor is essential. A fixed uncertainty added to the final scalar is a different error and does not gain this cubic decay. The uniform proof for arbitrary real $u$ needs no numerical square root; the propagated allowance remains conservative, and the exact rational fixtures check its location. The reverse proof uses different analytic state bounds, $9u^{3/2}$ and $(9/5)u^{3/2}$, with no numerical root needed in its theorem.

AI3's factored moments imply $|\Delta_3|\le4e$, $|\Delta_5|\le64e$, $|\Delta_7|\le832e$. Define exact rational constants
\begin{equation}
 B_*:=\frac23V^3+\frac8{15}V^5+\frac{52}{315}V^7,
 \qquad T_*:=\frac{91(6V)^9}{9!(1-6V)},
 \qquad A_*:=\frac{(6V)^9}{9!}.
 \label{r28:eq:ai4-retained-constants}
\end{equation}
The complete retained combination satisfies $|C_8|\le eB_*$ and its remaining tail is at most $eT_*$. Only fixed-source retained moment polynomials enter this statement; no physical error is canceled or differentiated. The ordinary individual allowance $A_*$ can remain constant in a positive denominator. Replacing the vanishing numerator allowance by a constant individual Taylor error would not establish this family for arbitrarily small $u$.

Nonnegative retained moments and alternating center signs give
\begin{equation}
 a_4\ge d_{\rm poly}:=\frac{zq_*^4}{100}-36V^3-\frac{(6V)^7}{7!}.
\end{equation}
The positive fifth-order contribution is dropped only for this lower bound. Let
\begin{equation}
 R_*:=\frac{36}{5}u_0^{3/2}+576\delta u_0^3+8D_6
       +48u_0^3M_0(1+3zM_0).
\end{equation}
The continuous bounds and all-band recurrences above imply $R_h\le R_*$ for both candidates. Exact rational evaluation proves
\begin{equation}
 d_{\rm poly}-A_*-R_*>d_*:=z/120>0,\qquad
 |y_5|<n_*:=z/32.
 \label{r28:eq:ai4-full-denominator}
\end{equation}
This bounds actual full scalars and the exported scalar-center denominator intervals, not just the retained matrix denominator.

\subsection{Explicit uniform gap and the independent constants}
Physical errors remain adversarially bounded:
\begin{equation}
 |y_5-qy_4|\le e(B_*+T_*)+(1+q)R_h
 \le e(B_*+T_*)+2R_h.
 \label{r28:eq:ai4-full-numerator}
\end{equation}
For $h=1,2$, put $(a_1,a_2)=(36/5,36/25)$ and $(b_1,b_2)=(576,1536/5)$, and define six exact contributions
\begin{align}
 c_h^{\rm state}&=\frac{2a_h\sqrt{u_0}}{d_*},&
 c_h^{\rm round}&=\frac{2b_h\delta u_0^2}{d_*},\notag\\
 c_h^{\rm spatial}&=\frac{32D_6}{d_*u_0},&
 c_h^{\rm average}&=\frac{96u_0^2M_0(1+3zM_0)}{d_*},\notag\\
 c_h^{\rm bias}&=\frac{hB_*}{d_*},&
 c_h^{\rm tail}&=\frac{hT_*}{d_*}.
 \label{r28:eq:ai4-radius-coefficients}
\end{align}
The band lower endpoint controls the spatial term; the upper endpoint controls the remaining physical terms. With $c_h$ their sum, every real $u$ in every band obeys
\begin{equation}
 \frac{y_5}{y_4}\in[q_h-c_hu,q_h+c_hu],\qquad
 \operatorname{gap}\ge\beta u,\qquad
 \beta:=1-c_1-c_2\approx0.9791915454711759>39/40.
 \label{r28:eq:ai4-uniform-result}
\end{equation}
The definition makes $\beta$ an explicit rational constant after the frozen substitutions; its full fraction is retained in the output. It is a sufficient coefficient, not an optimal physical margin or a measured signal-to-noise ratio.

The reverse producer obtains a different valid coefficient $\Gamma\approx0.975932912473311>39/40$ using its envelope and denominator $d>z/125$. Its compact formula is
\begin{equation}
 \Gamma=1-\frac{3(B_*+T_*)+2(P_1+P_2)}{d},
 \label{r28:eq:ai4-reverse-constant}
\end{equation}
where $P_h$ are its complete physical-radius coefficients, $R_h\le uP_h$. These include the $9,9/5$ state constants, the full spatial tail, and the proved growing-collar averaging bound. Its actual $d$ is the explicit rational lower bound from its scalar-center estimate, not the relaxed number $z/125$. The independent skeptic obtains a third sufficient coefficient near $0.9753467151692543$. These are independently derived certificates with different conservative choices; the common public conclusion is the weaker $39u/40$ floor.

In the forward bound the summed deficit is dominated by state, $0.020736$, and spatial, approximately $7.24545\times10^{-5}$. Averaging, correctly propagated rounding, retained bias and the full retained tail are also charged. Shortening duration reduces the spatial certificate strongly while increasing the relative state cost; the new schedule succeeds only because the combined bound remains below one.

\subsection{Vanishing scalar tolerance, exact fixtures and limits}
For an additional absolute error at most $\nu$ in each scalar, \eqref{r28:eq:ai4-full-denominator} gives a ratio change at most
\begin{equation}
 L\nu,\qquad L=\frac2{d_*}(1+n_*/d_*)=\frac{19}{2d_*},
 \qquad\nu\le d_*/2.
\end{equation}
The explicit forward choice
\begin{equation}
 \nu(u)=\frac{d_*u}{100}=\frac{zu}{12000}
 \label{r28:eq:ai4-scalar-tolerance}
\end{equation}
charges both candidates and leaves gap at least $(\beta-19/100)u>u/2$. At $u_0$ the allowance is about $8.33333\times10^{-36}$. It decreases to zero with $u$: no fixed positive apparatus error is covered uniformly. The four accepted allowances have distinct denominators and remaining margins:
\begin{center}\small
\begin{tabular}{@{}lll@{}}\toprule
Derivation & Additional error per scalar & Remaining gap\\\midrule
Forward & $d_*u/100$ & $(\beta-19/100)u>u/2$\\
Reverse, half margin & $\Gamma du/24$ & $\Gamma u/2$\\
Reverse, target margin & $(\Gamma-1/2)du/12$ & $u/2$\\
Skeptic, $D=z/128$ & $39Du/1280$ & $39u/80$\\\bottomrule
\end{tabular}
\end{center}
Each denominator reserve is proved in its own derivation. No mixed allowance/margin theorem is asserted, and the half-margin choices need not preserve $u/2$. Supplied preparation, measurement or calibration uncertainty must fit the applicable total scalar allowance.

Exactly ten physical parameter fixtures are evaluated: $j\in\{0,1,8,64,512\}$, each at $u=U_j$ and $u=3U_j/4$. They retain exact $q<1$, rational centers, complete errors and the proved face envelopes. At the two existing $j=0$ points only, the admitted exact $N_6=1808$ is compared with the envelope; this is not an extra cell. No new collar beyond six is computed. Large-band decimal displays that round a normalized gap to one do not erase its exact nonzero deficit.

Fixture success checks implementation. The theorem's continuous and all-integer quantifiers come from the geometry induction, continuous parameter inequalities and recurrences above. Separate algebraic controls reject finite-index extrapolation, wrong endpoint assignment, incomplete strip support, a fixed collar count, dropping $M^2$, deleting physical error terms, a wrong candidate clock, and confusing an inner arithmetic allowance with a final scalar floor. They do not supply extra physical observations or an automatic enlarged-graph Gram or omitted-channel theorem.

The shorter-duration strategy was proposed in the Feynman-method review after AI3; the advisor proposed its dyadic extension. These shared premises are attributed, while the producers and skeptic independently establish their consequences. The measured ratio contains no unknown $q$; candidate-dependent residuals and the time schedule are predictions for predeclared pairs. This is not a universal instrument setting secretly supplied an unknown $u$, or an inverse from fixed unrestricted data.

The common reference carrier and physical scale must be independently fixed. An unknown common phase need not preserve an imaginary-only ratio, and unsupplied timestamp or phase error requires a separate sensitivity estimate. There is no practical clock, apparatus, fixed-error robustness, homogeneous-model matching, global $q=1$ Hamiltonian, continuum construction or Yang--Mills mass-gap theorem. Scientific priority remains unverified.
